% =====================================================================================
% Chapter 5 · Price elasticity: hierarchical demand
%
% Built from notebooks/03_price_elasticity.ipynb. NOT ONE NUMBER IN THIS FILE IS TYPED BY
% HAND: every figure in the prose is a macro from build/macros.tex, generated by cmp.report
% from the executed notebook. The only literals below are (i) the constants of the
% data-generating model in §5.2, which are *definitions* rather than results, and (ii) the
% decision-rule thresholds (0.9 / 0.5) and the elasticity line at −1, which are *conventions
% and economics* stated in the text, not measurements.
% =====================================================================================
\chapter{Price elasticity: hierarchical demand}
\label{ch:pe}

\begin{quote}\small
\emph{Twelve regions, one price list, \nbThreeNObs{} region-weeks of sales --- and three regions
carrying only \nbThreeObsMin{} weeks apiece. Run a regression on those sixteen weeks and it
returns an elasticity of $\nbThreeWorstNopool$ where the planted truth is $\nbThreeWorstTrue$.
This is not an academic quantity: the elasticity \emph{is} the markup, by Lerner's rule, and the
sign of the pricing decision flips at $-1$. A hierarchical model with regional random slopes
repairs that region --- it is worth \pc{\nbThreeGainMean} of mean absolute error across
\nbThreeNseed{} fresh panels, though it \emph{loses} on \nbThreeNLosses{} of them --- and prices
\nbThreeNSet{} of the twelve regions with confidence. What it does not repair is that the firm
sets its own prices. On a panel where price responds to a demand shock the analyst cannot see,
the hierarchy's fleet elasticity lands within \nbThreeHierVsOls{} of naive OLS on the same data:
partial pooling shrinks noise, and endogeneity is bias. \Cref{sec:pe:wrong} watches the pricing
call reverse at $\kappa \approx \nbThreeKappaCross$, and hands the problem to an instrument.}
\end{quote}

% =====================================================================================
\section{The decision}
\label{sec:pe:decision}

A firm sells one product in \nbThreeNRegions{} regions and prices it region by region. The
question on the table is not academic: \emph{which regions should have their price moved, in
which direction, and by how much?} A tourist market and a price-sensitive suburb should not carry
the same ticket, and the quantity that decides the difference has a name.

The \textbf{price elasticity of demand} is the percentage change in quantity sold for a one per
cent change in price. It is almost always negative --- raise the price, sell less --- and an
elasticity of $-1.5$ reads: cut price by one per cent, sell one and a half per cent more. The
firm needs one such number per region, and it needs it to be causal: what happens to demand
\emph{when we move the price}, not what happened to demand on the weeks a manager chose to
discount.

\subsection*{Why the number matters, and why the sign flips at \texorpdfstring{$-1$}{-1}}

Write demand as a constant-elasticity curve $Q(P) = A P^{\beta}$ and revenue as
$R(P) = P\,Q(P) = A P^{1+\beta}$. Then
%
\begin{equation}
  \frac{dR}{dP} \;=\; (1+\beta)\, A\, P^{\beta}
  \qquad\Longrightarrow\qquad
  \frac{dR}{dP} > 0 \iff \beta > -1 .
  \label{eq:pe:revenue}
\end{equation}
%
\Cref{eq:pe:revenue} is the whole reason $-1$ is not an arbitrary threshold. Above it (demand
\emph{inelastic}) a price \emph{rise} increases revenue and, since fewer units are sold, cuts cost
too --- the firm is under-charging. Below it (demand \emph{elastic}) a price \emph{cut} increases
revenue. A recommendation that gets $\beta$ wrong by a tenth of a point is a rounding error; a
recommendation that gets it wrong \emph{across} $-1$ inverts the instruction sent to the region
manager. That is the asymmetry a pricing model must be built around, and it is why
\cref{sec:pe:euro}'s decision rule is written as $\Prob(\beta_r < -1)$ rather than as ``is the
estimate statistically significant?''. Significance is a statement about zero. The business does
not care about zero. It cares about $-1$.

The profit optimum sharpens the point. With constant marginal cost $c$, weekly profit is
$\pi(P) = (P-c)\,A P^{\beta}$, and setting $d\pi/dP = 0$ gives
%
\begin{equation}
  P^\star \;=\; \frac{\beta}{\beta+1}\,c
  \qquad\Longleftrightarrow\qquad
  \frac{P^\star - c}{P^\star} \;=\; -\frac{1}{\beta},
  \label{eq:pe:lerner}
\end{equation}
%
the \textbf{Lerner rule} \citep{lerner1934}: \emph{the optimal markup over cost is the inverse
elasticity}. An elasticity of $-2$ prices at twice cost, $-1.5$ at three times cost. The
elasticity is not an input to the pricing decision --- under this model it very nearly \emph{is}
the pricing decision. And \cref{eq:pe:lerner} is undefined for $\beta \ge -1$: the denominator
$\beta + 1$ crosses zero, $P^\star$ blows up, and the model politely announces that it has left
its domain. \Cref{sec:pe:euro} makes that instability do useful work.

\subsection*{What makes it hard}

Three things, and they are the chapter.

\textbf{First, the regions are not the same size.} Nine of the twelve carry \nbThreeObsMax{} weeks
of history; three carry \nbThreeObsMin{}. A separate regression on sixteen weeks returns an
elasticity that is mostly sampling noise, and acting on it is not boldness but arithmetic error.
The classical toolbox offers exactly two ways out --- pretend the regions are identical, or trust
each region's own numbers --- and both are wrong in opposite directions. This is the same
bias--variance bracket as Chapter~4, now with a coefficient rather than a mean, and this
chapter's third option is a hierarchical model that learns \emph{how much} to pool.

\textbf{Second, the model is chosen for interpretability, not convenience.} Take logs of both
sides of $Q = A P^{\beta}$ and the demand curve becomes linear:
$\log Q = \log A + \beta \log P$. The slope on log-price \emph{is} the elasticity --- no
back-transformation, no marginal effect to compute, no delta method. One line of calculus says
why: elasticity is defined as $\epsilon(P) \equiv \frac{dQ}{dP}\cdot\frac{P}{Q}$, and for
$Q = A P^{\beta}$,
%
\begin{equation}
  \epsilon(P) \;=\; \beta A P^{\beta - 1} \cdot \frac{P}{A P^{\beta}} \;=\; \beta
  \qquad \text{at every price.}
  \label{eq:pe:elast}
\end{equation}
%
For any \emph{other} demand shape the elasticity moves with the price level, so $\beta$ is best
read as a local approximation around the prices actually observed --- a caveat that returns with
teeth in \cref{sec:pe:euro}, where every recommended price sits at or beyond the edge of that
range.

\textbf{Third, and hardest: price is not assigned by a coin.} The firm sets it. Managers
historically cut prices when they expected soft weeks, which means low prices and low demand
coincide for reasons that have nothing to do with the price cut. This is \emph{endogeneity}, and
it is the failure mode that no amount of clever averaging repairs. \Cref{sec:pe:wrong} shows the
hierarchical posterior inheriting the bias to within \nbThreeHierVsOls{} of naive OLS, and shows
the pricing recommendation reversing outright. The cure is an instrument, and this chapter hands
the problem to Chapter~13.

% =====================================================================================
\section{The data-generating model}
\label{sec:pe:dgm}

A causal method cannot be graded on real data, because on real data the causal slope is exactly
what is missing. So the method is graded first against a world whose elasticities are known by
construction.

The simulator produces a region-week panel: $R = \nbThreeNRegions$ regions, indexed
$r = 1, \dots, R$, over \nbThreeNWeeks{} weeks. Each region carries its own true
elasticity, drawn once and then fixed; the season $s_t$ and the trend $g_t$ are shared; the
competitor's price $C_{rt}$ and a demand shock $u_{rt}$ are drawn afresh each week:
%
\begin{align}
  \beta_r &\sim \mathcal N\!\left(-1.4,\ 0.35^{2}\right), \label{eq:pe:beta}\\
  s_t &= 0.3 \sin\!\Big(\tfrac{2\pi t}{26}\Big), \qquad g_t = 0.01\,t, \qquad
        C_{rt} \sim \mathcal N\!\left(20,\ 2^{2}\right), \qquad
        u_{rt} \sim \mathcal N(0, 1), \label{eq:pe:latent}\\
  \log P_{rt} &= \log 15 \;-\; 0.05\,\frac{C_{rt}}{20} \;+\; \kappa\, u_{rt} \;+\; \nu_{rt},
        \qquad \nu_{rt} \sim \mathcal N\!\left(0,\ 0.05^{2}\right), \label{eq:pe:price}\\
  \log Q_{rt} &= 2.5 \;+\; \beta_r \log P_{rt} \;+\; 0.4\,\frac{C_{rt} - 20}{20} \;+\; s_t \;+\; g_t
        \;+\; u_{rt} \;+\; \varepsilon_{rt},
        \qquad \varepsilon_{rt} \sim \mathcal N\!\left(0,\ 0.08^{2}\right). \label{eq:pe:demand}
\end{align}
%
\Cref{eq:pe:beta} plants the twelve elasticities. \Cref{eq:pe:price} is the firm's price-setting
rule, and \cref{eq:pe:demand} is demand. Read the two together, because the whole chapter is in
their overlap.

The competitor's price $C_{rt}$ appears in \emph{both} equations, and that is precisely what makes
it a \textbf{confounder}: it lowers our price (the $-0.05\,C/20$ term) and raises our demand (the
$+0.4$ term), so an analyst who ignores it sees cheap weeks coinciding with strong weeks and
credits the price cut for demand the competitor delivered. Season and trend appear only in
\cref{eq:pe:demand}: they move demand and not price, so they are \textbf{precision covariates} ---
worth controlling for, but not because they bias anything.

$\kappa$ is the \textbf{endogeneity switch}. Through \crefrange{sec:pe:classical}{sec:pe:euro} the
generator runs at $\kappa = 0$ and zeroes the shock entirely ($u_{rt} \equiv 0$), so recovery is
clean and every estimator can be graded on a level field. \Cref{sec:pe:wrong} sets
$\kappa = \nbThreeKappaDemo$, and the \emph{same} shock $u_{rt}$ then enters both the price
equation and the demand equation --- managers responding to the very thing that moves demand. That
is the simultaneity that no pooling scheme can undo.

Two features of the planted world set the difficulty. The fleet-mean elasticity is
$\nbThreeTrueFleet$ and the between-region spread is $\nbThreeTrueTau$, so the regions genuinely
differ --- a single national price is leaving money somewhere. And of the \nbThreeNRegions{}
planted elasticities, \nbThreeNTrueElastic{} lie below $-1$ and exactly \nbThreeNTrueInelastic{}
lies above it (\nbThreeInelasticRegion{}, at $\nbThreeInelasticBeta$). That one region is the trap:
any estimator that certifies all twelve as elastic has committed a \emph{wrong} action, not merely
a cautious one. Finally, three regions are deliberately thinned to \nbThreeObsMin{} weeks, leaving
\nbThreeNObs{} region-weeks in total, so that shrinkage has something to do.

\Cref{fig:pe:dag} draws the graph these equations imply.

\input{build/floats/nb03_dag}

% =====================================================================================
\section{Identification}
\label{sec:pe:ident}

\subsection*{Identification is not estimation}

The distinction governs the rest of the chapter, so it is worth stating flatly.
\emph{Identification} asks whether the causal quantity could be recovered from this data even with
infinitely many weeks --- a question about assumptions, settled before any model is fitted.
\emph{Estimation} is the statistics done once identification has been granted. Everything this
chapter does with pooling --- complete, none, partial --- is an \textbf{estimation} device. It buys
variance reduction. It never buys identification, and \cref{sec:pe:wrong} is the demonstration.

\subsection*{The estimand}

Region $r$'s elasticity is a causal slope: what demand does when \emph{we} intervene on price, not
when the market moves it for us. In the do-calculus of \citet{pearl2009},
%
\begin{equation}
  \beta_r \;=\;
  \frac{\partial\, \E\!\big[\log Q \,\big|\, \operatorname{do}(\log P = p),\ r\big]}{\partial p},
  \label{eq:pe:estimand}
\end{equation}
%
where $\operatorname{do}(\cdot)$ marks an intervention, as opposed to passively observing whatever
prices managers happened to choose. The equivalent potential-outcome statement
\citep{rubin1974, imbens2015} is that $\beta_r$ is the slope of region $r$'s potential-outcome
curve $p \mapsto \log Q_{r}(p)$, of which exactly one point is observed each week. Every estimator
in this chapter --- pooled OLS, twelve separate regressions, the hierarchical posterior, and 2SLS
--- targets \cref{eq:pe:estimand}. Keeping the estimand fixed is what makes \cref{sec:pe:compare} a
comparison rather than a change of subject.

\subsection*{The assumptions}

The regression slope equals the causal slope of \cref{eq:pe:estimand} under four assumptions. Each
is stated, and each is either tested later or explicitly declared untestable --- an assumption with
no accompanying test is a wish.

\begin{assumption}[Conditional price exogeneity]\label{as:pe:exog}
Once the observed controls $Z_{rt}$ (competitor price, season, trend) and the region are held
fixed, price is as good as randomly assigned with respect to the demand shock:
%
\[
  \varepsilon_{rt} \;\perp\; \log P_{rt} \;\big|\; Z_{rt},\, r .
\]
%
This is the continuous-treatment form of unconfoundedness. It holds by construction while
$\kappa = 0$, and it is \textbf{untestable from the data alone}: no diagnostic in this chapter can
certify it, and \cref{fig:pe:ppc} shows a model that violates it sailing through every check a real
analyst can run. It is the assumption \cref{sec:pe:wrong} breaks on purpose.
\end{assumption}

\begin{assumption}[Price variation]\label{as:pe:variation}
Genuine price movement survives the controls:
%
\[
  \Var(\log P_{rt} \mid Z_{rt},\, r) \;>\; 0 .
\]
%
No variation, no slope. This is positivity, wearing continuous-treatment clothes. It is
testable, and tested: the within-region standard deviation of log-price runs from
\nbThreeLpSdMin{} to
\nbThreeLpSdMax{} (median \nbThreeLpSdMed), which is the quasi-experimental wiggle
$\nu_{rt} \sim \mathcal N(0, 0.05^2)$ of \cref{eq:pe:price} showing up in the data. In a real firm
this is the assumption most often \emph{silently} violated: uniform corporate pricing or a rigid
cost-plus rule leaves $\Var(\log P \mid Z) \approx 0$, and the ``elasticity'' estimated from it is
noise wearing a number's clothes.
\end{assumption}

\begin{assumption}[No dynamics, no cross-region spillover]\label{as:pe:sutva}
This week's price affects this week's demand only --- no stockpiling ahead of a price rise, no
reference-price memory --- and one region's price does not move another region's demand. This is
SUTVA \citep{imbens2015} applied to a pricing panel. It holds by construction here. On real scanner
data it is the first thing to probe, with a lagged-price term: pantry-loading masquerades as high
short-run elasticity and evaporates at longer horizons.
\end{assumption}

\begin{assumption}[Constant elasticity]\label{as:pe:form}
Within a region, the demand curve really is log-log, so a single $\beta_r$ governs every price. In
a parametric model the functional form is part of the identification bundle, not a convenience: if
the true curve bends, the ``constant'' $\beta_r$ is a price-range-specific average. Tested in
sample in \cref{sec:pe:valid} with a curvature term --- and no in-sample test reaches prices never
charged, which is exactly where \cref{eq:pe:lerner} wants to send the firm.
\end{assumption}

\subsection*{What it costs when \texorpdfstring{\cref{as:pe:exog}}{Assumption 5.1} fails}

Not hand-waving --- two lines of algebra. Let an unobserved shock $u_{rt}$, with variance
$\sigma_u^2$, enter demand with coefficient $1$ and leak into price with coefficient $\kappa$
(managers cut price when they expect a soft week). Regress $\log Q$ on $\log P$ after partialling
out $Z$ and the region effects. The classical omitted-variable-bias formula gives
%
\begin{equation}
  \hat\beta \;\xrightarrow{\ p\ }\; \beta \;+\;
  \frac{\operatorname{Cov}(\log P,\, u)}{\Var(\log P \mid Z)}
  \;=\; \beta \;+\; \frac{\kappa\,\sigma_u^{2}}{\kappa^{2}\sigma_u^{2} + \sigma_\nu^{2}},
  \label{eq:pe:ovb}
\end{equation}
%
reading the covariance and the variance straight off \cref{eq:pe:price}. Three things follow, and
all three are measured in \cref{sec:pe:wrong}.

The bias term is \textbf{positive}, and since $\beta < 0$ it drags $\hat\beta$ \emph{toward zero}:
demand looks less price-sensitive than it is, and the firm is told to hold a price it should have
cut. The bias does \textbf{not shrink with more data} --- $\xrightarrow{p}$ reads ``converges in
probability to'', and it converges to the wrong number. And the bias is \textbf{largest when the
honest price variation is smallest}: the denominator carries $\sigma_\nu^{2}$, which is $0.05^2$
here, so even a faint $\kappa$ swamps it. Push $\kappa$ far enough and the fitted demand curve
slopes \emph{upward} --- high prices coinciding with high demand, because both are chasing the same
shock. \Cref{fig:pe:kappa} plots \cref{eq:pe:ovb} against the refits and they agree across the whole
sweep, perverse region included.

This is the mechanism that makes price a special case in demand estimation, and it is why the
industrial-organization literature reaches for instruments as a matter of course rather than as a
robustness check \citep{berry1994, blp1995, nevo2000}. \Cref{as:pe:exog} is the one assumption in
this chapter that no amount of estimation skill can rescue.

% =====================================================================================
\section{The classical baseline}
\label{sec:pe:classical}

Before any sampler, the discipline of this book is to ask what a competent analyst produces
without one. The answer is not a different analysis: it is the \emph{same} estimand
(\cref{eq:pe:estimand}), identified by the \emph{same} four assumptions, estimated with the
simplest tool that is still correct. And here that tool is almost embarrassingly simple, because
the estimand \emph{is} a regression coefficient:
%
\begin{equation}
  \log Q_{rt} \;=\; \alpha_r \;+\; \beta\, \log P_{rt} \;+\; \gamma^{\top} Z_{rt}
                \;+\; \varepsilon_{rt},
  \label{eq:pe:pooled}
\end{equation}
%
with region intercepts $\alpha_r$ (twelve dummies) and $Z$ the adjustment set of
\cref{sec:pe:ident}. The coefficient on $\log P$ is the elasticity. No sampler is required to read
it.

\subsection*{The covariance choice is the estimate's other half}

The point estimate of \cref{eq:pe:pooled} is $\nbThreeFleetOls$ --- $\nbThreeFleetErr$ from the
planted fleet mean of $\nbThreeTrueFleet$, which is a good number. But an estimate without an
interval is not a result, and the interval depends entirely on what one is willing to assume about
the errors. \Cref{tab:pe:covariance} fits \cref{eq:pe:pooled} three times, changing nothing but
that assumption.

\input{build/tables/nb03_covariance}

The \nbThreeNObs{} region-weeks are \emph{not} independent observations, and for two distinct
reasons, of which only the second is the textbook one.

\textbf{Reason one: the pooled model is deliberately wrong about slopes.} \Cref{eq:pe:beta} planted
$\beta_r \sim \mathcal N(-1.4,\, 0.35^2)$ --- the regions genuinely differ. Force a single shared
$\beta$ on them, as \cref{eq:pe:pooled} does, and each region's residual inherits a term
$(\beta_r - \beta)\log P_{rt}$ that carries \emph{the same sign, week after week, for that region}.
That is within-region correlation manufactured by the specification itself, and it exists even
though the simulator's weekly shocks are iid by construction.

\textbf{Reason two: real panels add serially correlated local shocks} on top --- a soft month tends
to be followed by another soft month.

Neither is visible to the iid or HC1 formulas, which only ever inspect one row at a time. HC1
\citep{white1980, mackinnon1985} is robust to heteroskedasticity and to nothing else. Clustering on
region \citep{liang1986} lets the errors be arbitrarily correlated \emph{within} a region and asks
only for independence \emph{across} regions, so the \nbThreeNClusters{} regions --- not the
\nbThreeNObs{} rows --- become the unit of inference. This is the correction of
\citet{bertrand2004}, whose finding is that ignoring within-panel correlation can leave standard
errors wrong by a factor of five or more.

The point estimate is identical to the last decimal under all three; only the story about its
precision changes, and it changes a lot. Clustering multiplies the standard error by
\nbThreeSeRatioIid$\times$ over iid and \nbThreeSeRatioHc$\times$ over HC1. Two lessons, and the
second is the one that is usually missed.

\begin{enumerate}[leftmargin=*]
  \item \textbf{The iid and HC1 intervals are not conservative or anti-conservative as a rule.}
        They are simply \emph{wrong in a direction you cannot predict} without doing the
        clustering. There is no safe default here to fall back on.
  \item \textbf{The widening is itself a diagnosis.} What the clusters found was not serial
        correlation --- the simulator planted none --- but reason one: the pooled model's own slope
        misspecification. The classical machinery is \emph{telling us} that one $\beta$ is not
        enough for twelve regions. It simply has no way to act on the message. That is exactly the
        message \cref{sec:pe:bayes} acts on.
\end{enumerate}

\subsection*{And even the best interval here is the optimistic one}

One caveat belongs \emph{before} the number is quoted, not in a footnote after it. There are
\nbThreeNClusters{} clusters. Cluster-robust standard errors are a large-$G$ approximation, and at
$G = \nbThreeNClusters$ they are themselves \emph{downward}-biased: the usual rule of thumb wants
thirty to fifty clusters, and below that the recommended repair is a wild cluster bootstrap
\citep{cameron2008, cameron2015}. So the clustered interval of \cref{tab:pe:covariance},
$[\nbThreeFleetLo,\ \nbThreeFleetHi]$, is the best interval the classical toolbox has in this
problem --- and it is still, mildly, an optimistic one. No covariance option fixes that. The honest
fix is to \emph{model} the region-level variation, which is what \cref{sec:pe:bayes} does. This is
the register in which a covariance estimator should be chosen: pick the one that makes the fewest
false promises, then say out loud which promises it is still making.

\subsection*{The control ladder: bias and noise are different failures}

\Cref{sec:pe:ident} claimed that competitor price is the one genuine confounder and that season and
trend are mere precision covariates. That is a claim about the data-generating model, and it is
testable by dropping the controls one block at a time. But a single panel cannot tell \emph{bias}
from \emph{bad luck}, so the ladder is re-run on \nbThreeNseed{} fresh panels and the \emph{average
error} reported (\cref{tab:pe:ladder}).

\input{build/tables/nb03_ladder}

Dropping the confounder moves the estimate by $\nbThreeBiasNocomp$ on average --- small, but
\emph{systematic}, and pointing away from zero, so demand looks more price-sensitive than it is.
The sign is exactly what \crefrange{eq:pe:price}{eq:pe:demand} dictate. Unlike noise, this error
does not shrink with more data. It is small here only because the competitor passthrough is weak by
construction; with a strong one the same mechanism is ruinous.

Dropping season and trend is a different animal entirely. The mean error stays near zero
($\nbThreeBiasNoseason$) --- no bias, exactly as the graph advertises --- but the typical standard
error widens by \nbThreeSeWiden$\times$ (from \nbThreeSeLadderFull{} to
\nbThreeSeLadderNoseason), and the seed-to-seed scatter inflates to match. On the \emph{working}
panel it happened to throw the point estimate a long way, which is noise, and which is precisely
what one sample can never tell you and \nbThreeNseed{} panels can. The general lesson outlives this
notebook: when the treatment barely moves --- and here $\log P$ has a standard deviation of about
$0.05$ --- anything that explains a large share of the \emph{outcome} is worth controlling for even
when it confounds nothing at all.

\subsection*{Where the classical arm runs out of road}

It landed the fleet number. The business asked for twelve.

And here the classical toolbox has exactly two options, neither of which is good.
\textbf{Complete pooling} is \cref{eq:pe:pooled} itself: one $\beta$, quoted twelve times. It is
rock-steady and it is biased for every region that genuinely differs --- which, by
\cref{eq:pe:beta}, is all of them. \textbf{No pooling} runs twelve separate regressions, each on
its own region's weeks. It is unbiased and, on a \nbThreeObsMin-week region, close to useless: on
\nbThreeWorstRegion{} it returns $\nbThreeWorstNopool$ against a planted $\nbThreeWorstTrue$
(\cref{tab:pe:regions}), an error of nearly a whole elasticity point in a quantity whose decision
line is at $-1$.

There is no third option in the classical toolbox.\footnote{A frequentist mixed model would get
close --- it is, after all, the same likelihood. The point is that plain OLS, the tool actually
reached for, cannot.} The fleet elasticity under no pooling averages $\nbThreeNopoolFleet$ and
under complete pooling is $\nbThreeCompleteFleet$; the disagreement between them, region by region,
is the entire subject of the next two sections. \Cref{sec:pe:bayes} supplies the third option, and
\cref{sec:pe:compare} settles the account between them at a bar that has been matched honestly.

% =====================================================================================
\section{The Bayesian model}
\label{sec:pe:bayes}

The hierarchical move is one line. Instead of asserting that the regions are identical (complete
pooling) or that they have nothing to say to each other (no pooling), assert that they are
\emph{drawn from a common distribution} whose spread is estimated from the data:
%
\begin{align}
  \log Q_i &\;\sim\; \mathcal N\Big(\alpha_{r[i]} + \beta_{r[i]} \log P_i
              + \gamma^{\top} z_i,\ \ \sigma^{2}\Big), \label{eq:pe:lik}\\[2pt]
  \beta_r &\;\sim\; \mathcal N\!\left(\mu_\beta,\ \tau_\beta^{2}\right), \qquad
  \alpha_r \;\sim\; \mathcal N\!\left(0,\ 5^{2}\right), \qquad
  \gamma \;\sim\; \mathcal N\!\left(0,\ 2^{2}\right), \label{eq:pe:slopes}\\[2pt]
  \mu_\beta &\;\sim\; \mathcal N\!\left(-1,\ 1^{2}\right), \qquad
  \tau_\beta \;\sim\; \text{HalfNormal}(1), \qquad
  \sigma \;\sim\; \text{HalfNormal}(1), \label{eq:pe:priors}
\end{align}
%
with $i$ indexing region-weeks, $r[i]$ region $i$'s region, and $z_i$ the standardized controls ---
the \emph{same} adjustment set the classical arm used, so that nothing in the comparison turns on a
different specification. \Cref{eq:pe:slopes}'s middle term is the whole idea: $\mu_\beta$ is the
typical elasticity, $\tau_\beta$ is how much regions genuinely differ, and both are learned.

\subsection*{Why these priors}

The priors are weak, and the chapter's honesty depends on saying so rather than hoping nobody asks.
$\mu_\beta \sim \mathcal N(-1, 1)$ puts roughly 95\,\% of its mass on $(-3, +1)$: it spans
everything from mildly inelastic demand to an elasticity of $-3$, which is more price-sensitive
than almost any retail elasticity ever measured --- published meta-analyses of scanner-data
elasticities centre well inside that band \citep{tellis1988, bijmolt2005} --- while gently
centring on the economically plausible region. With \nbThreeNObs{} region-weeks the likelihood
dominates it easily, and the posterior fleet mean, $\nbThreeMuBeta$, lands within
\nbThreeMuVsOls{} of the pooled OLS estimate that used no prior at all.

$\tau_\beta \sim \text{HalfNormal}(1)$ is the load-bearing one, and it is worth seeing what it
does. As $\tau_\beta \to 0$ the model collapses to \textbf{complete pooling} (all regions
identical); as $\tau_\beta$ grows large it approaches \textbf{no pooling} (every region on its
own). The two classical baselines of \cref{sec:pe:classical} are the \emph{limiting cases} of the
model actually being fitted, and where the fit lands between them is \emph{learned, not assumed}.
The posterior puts $\tau_\beta$ at \nbThreeTauBeta{} (90\,\% credible interval
$[\nbThreeTauBetaLo,\ \nbThreeTauBetaHi]$) against a planted \nbThreeTrueTau. The data chose the
pooling regime, and it chose neither extreme.

\subsection*{Shrinkage, and what it is worth}

The consequence of \cref{eq:pe:slopes} is that each region's posterior elasticity is a
precision-weighted compromise between its own weeks and the fleet mean. The pull toward the fleet
--- the \textbf{shrinkage factor} --- is approximately
%
\begin{equation}
  w_r \;=\; \frac{se_r^{2}}{\tau_\beta^{2} + se_r^{2}}
  \;\approx\; \frac{\sigma^{2}/n_r}{\tau_\beta^{2} + \sigma^{2}/n_r},
  \qquad
  \hat\beta_r^{\,\text{partial}} \;\approx\; (1 - w_r)\,\hat\beta_r^{\,\text{no pool}}
                                   \;+\; w_r\, \mu_\beta ,
  \label{eq:pe:shrink}
\end{equation}
%
where $se_r$ is region $r$'s no-pooling slope standard error. Read it: a region with plenty of
weeks has a small $se_r$, hence a small $w_r$, and barely moves. A region with sixteen weeks has a
large $se_r$, hence a large $w_r$, and leans hard on the fleet. The pull is not a fudge factor ---
it is the exact weight that minimizes expected squared error if \cref{eq:pe:slopes} is true, and it
is the same estimator, arrived at from a different direction, that \citet{stein1956} and
\citet{james1961} showed dominates the twelve separate maximum-likelihood estimates, and that
\citet{efron1977} popularized as ``borrowing strength''. \Citet{gelmanhill2007} is the standard
modern treatment, and this chapter is one of its examples with money attached.
\Cref{fig:pe:shrink} verifies \cref{eq:pe:shrink} against the pull the sampler actually applied.

\subsection*{The sampler}

Inference is by the No-U-Turn sampler \citep{hoffman2014} in PyMC \citep{salvatier2016}:
\nbThreeNDraws{} post-warmup draws, $\hat R = \nbThreeRhat$, a minimum bulk effective sample size
of \nbThreeEss{} \citep{vehtari2021}, and \nbThreeDivergences{} divergent transitions. Divergences
are the diagnostic that matters most in a hierarchy, because the centred parameterization of
\cref{eq:pe:slopes} has a funnel geometry --- when $\tau_\beta$ is small the $\beta_r$ are squeezed
into a narrow neck the sampler cannot explore with a step size tuned for the wide part
\citep{betancourt2015}. Here $\tau_\beta$ is comfortably away from zero and the centred
parameterization samples cleanly. Were it not, the repair is a non-centred reparameterization, not
a stronger prior. The chains have converged, and nothing that follows is a sampler artefact ---
including, importantly, the \nbThreeNseed{} refits of \cref{sec:pe:compare}, across which the worst
$\hat R$ was \nbThreeRefitRhat{} and the divergence count was \nbThreeRefitDiv{} \emph{in total} ---
about one per refit, against \nbThreeNDraws{} draws apiece, which is immaterial.

% =====================================================================================
\section{Diagnostics and validation}
\label{sec:pe:valid}

The discipline is to run the checks a \emph{real} analyst can run before the ones only a simulator
affords, and to keep firmly apart what each of them certifies.

\subsection*{The checks a real analyst can run}

\Cref{fig:pe:ppc} carries three. The posterior predictive check simulates replicated datasets from
the fitted model and overlays them on the observed log-demand: replicate mean \nbThreePpcMean{}
against an observed \nbThreeObsMean, replicate standard deviation \nbThreePpcSd{} against
\nbThreeObsSd. Residuals plotted against log-price probe \cref{as:pe:form}, because curvature
there is exactly what a violated constant-elasticity assumption looks like. Residuals plotted
against week probe \cref{as:pe:sutva}: leftover structure means missed seasonality or dynamics.

\input{build/floats/nb03_ppc}

\Cref{as:pe:form} also gets a formal probe: add $(\log P - \overline{\log P})^2$ to a pooled fit
with the same controls and ask whether its coefficient is distinguishable from zero. It is not
($\nbThreeCurvCoef$, standard error \nbThreeCurvSe, $t = \nbThreeCurvT$) --- constant elasticity
survives in sample. But read the \emph{width} of that standard error, because it is the more
informative half of the result: with only about $\pm 5\,\%$ of price variation to work with, this
test has almost no power to detect curvature far from current prices. Any $P^\star$ outside the
observed band is extrapolation, not measurement, and \cref{sec:pe:euro} will find that every
$P^\star$ is.

\begin{remark}[What a clean posterior predictive check does \emph{not} certify]
A clean PPC certifies \textbf{fit}, not \textbf{causality}. The endogenous model of
\cref{sec:pe:wrong} --- wrong by nearly a full elasticity point, and wrong in the direction that
reverses the pricing recommendation --- passes every one of these checks, because confounding is
invisible in the distribution of the outcome alone. PPC failures tell you the model is wrong; PPC
successes do not tell you it is right about the causal question. That was \cref{as:pe:exog}, and
\cref{as:pe:exog} is untestable.
\end{remark}

\subsection*{The checks only a simulator affords}

\input{build/floats/nb03_shrink}

\Cref{fig:pe:shrink} grades the posterior against the truth. The recovery panel puts the twelve
posterior elasticities on the $45^\circ$ line with intervals that are visibly wider where the data
is thin. The shrinkage panel shows \cref{eq:pe:shrink} in action: the three thin regions are yanked
toward the fleet, the data-rich ones barely move. And the third panel closes the loop on the
formula --- predicted pull \nbThreeShrinkPredThin{} against realized \nbThreeShrinkRealThin{} for
the thin regions, \nbThreeShrinkPredRich{} against \nbThreeShrinkRealRich{} for the rich ones. The
formula the reader was asked to trust in \cref{eq:pe:shrink} is the formula the sampler obeyed.

\subsection*{Calibration, and an honest look at a below-nominal number}

The 90\,\% credible intervals cover the twelve planted elasticities \pc{\nbThreeCover} of the time
--- \nbThreeCoverHits{} of \nbThreeNRegions. That is below nominal, and it should not be waved
away. Two checks, and they say different things.

\emph{Where do the misses live?} \nbThreeCoverThin{} of the three thin regions are covered and
\nbThreeCoverRich{} of the nine rich ones. This is the known price of shrinkage, and it is
structural rather than a bug: shrinkage intervals are calibrated for the \emph{ensemble}, not for
each pre-shrinkage slope, so they run too narrow precisely for regions whose true elasticity sits
far from the fleet mean and for the hard-shrunk thin ones.

\emph{Is the shortfall even signal?} With only \nbThreeNRegions{} regions, a perfectly calibrated
90\,\% procedure still misses a couple regularly. The binomial probability that a truly-90\,\%
method covers \nbThreeCoverHits{} or fewer of twelve is \nbThreePAsbad. Read that number before
reading anything into the coverage.

The trade is real and worth naming, because it is the trade the chapter's decision rule is built
around: partial pooling buys the \emph{best point estimates} (\cref{sec:pe:compare}) at the cost of
slightly overconfident \emph{per-region intervals}. That is one more reason \cref{sec:pe:euro}
leans on $\Prob(\beta_r < -1)$ --- a probability computed across the whole posterior --- rather
than on interval endpoints.

% =====================================================================================
\section{Point estimate versus posterior}
\label{sec:pe:compare}

Three arms, one estimand (\cref{eq:pe:estimand}), one set of assumptions, one adjustment set. They
differ \emph{only} in how the slope is pooled. \Cref{tab:pe:regions} puts them side by side, region
by region, against the planted truth.

\input{build/tables/nb03_regions}

\subsection*{On the fleet number, they agree --- say it plainly}

The pooled OLS of \cref{sec:pe:classical} says $\nbThreeFleetOls$. The posterior $\mu_\beta$ says
$\nbThreeMuBeta$, with a 90\,\% credible interval of
$[\nbThreeMuBetaLo,\ \nbThreeMuBetaHi]$. They differ by \nbThreeMuVsOls, and both sit within a hair
of the planted $\nbThreeTrueFleet$. This is not a disappointment; it is the lesson, and it is the
book's thesis arriving on schedule. \textbf{The causal work was done by the identification, not by
the machinery.} With \nbThreeNObs{} region-weeks and weak priors, a Bayesian hierarchy is not going
to move a well-identified pooled slope, and it does not --- the Bernstein--von Mises phenomenon
\citep{gelman2013, vandervaart1998}, doing what it always does. If the reason for reaching for a
sampler was a better headline number, this is the correction.

\subsection*{On the twelve numbers, partial pooling wins --- and the margin is not what one panel says}

Now the estimand the business actually asked for. \Cref{tab:pe:bakeoff} scores all three arms by
mean absolute error against the twelve planted elasticities --- first on the working panel, then on
\nbThreeNseed{} fresh ones, each thinned identically and each given a full hierarchical refit at
the same sampler settings.

\input{build/tables/nb03_bakeoff}
\input{build/floats/nb03_multiseed}

The working panel says partial pooling beats the better classical arm by \pc{\nbThreeGainOne}. That
number is an anecdote, and this chapter has already spent a section refusing to read bias off one
draw (\cref{tab:pe:ladder}); the headline gets no exemption. Over \nbThreeNseed{} panels:

\begin{enumerate}[leftmargin=*]
  \item \textbf{The average margin is \pc{\nbThreeGainMean}, not \pc{\nbThreeGainOne}}, with a
        standard deviation of \nbThreeGainSd\,pp. Across panels the edge runs from
        \pc{\nbThreeGainMin} --- on its worst panel partial pooling is \emph{worse}, and by a wide
        margin --- to \pc{\nbThreeGainMax}. The working panel sits at the \nbThreePctileOne th
        percentile of that spread: a flattering draw, not a representative one.
  \item \textbf{It loses on \nbThreeNLosses{} of \nbThreeNseed{} panels}, winning
        \nbThreeNWins{}. And the comparator is deliberately unhelpful to the Bayesian arm: it is
        the \emph{oracle-best} classical arm, re-chosen panel by panel against a planted truth no
        analyst would have. Partial pooling has to beat whichever extreme happened to get lucky.
        Head to head it beats no pooling on \nbThreeWinsVsNopool{} panels and complete pooling on
        \nbThreeWinsVsComplete.
  \item \textbf{Which classical arm is the one to beat is itself a one-seed artefact.} On the
        working panel, complete pooling (\nbThreeMaeCompleteOne) looks better than no pooling
        (\nbThreeMaeNopoolOne) --- but only because \nbThreeWorstRegion's sixteen weeks produced one
        spectacular miss that drags the no-pooling average up single-handed. Averaged over panels
        the ranking \textbf{reverses}: no pooling \nbThreeMaeNopool, complete pooling
        \nbThreeMaeComplete. With \nbThreeObsMax{} weeks in nine of twelve regions, ``trust each
        region's own data'' is usually the better classical bet, and complete pooling's
        identical-for-everyone bias is the steadier drag. The blow-up in the table is the anecdote;
        the mean is the fact.
\end{enumerate}

So the claim survives contact with \nbThreeNseed{} panels, and it lands quieter than one panel
would let anyone say it. Partial pooling has the lowest mean error (\nbThreeMaePartial{} against
\nbThreeMaeNopool{} and \nbThreeMaeComplete) \emph{and} the smallest seed-to-seed spread
(\nbThreeMaePartialSd{} against \nbThreeMaeNopoolSd{} and \nbThreeMaeCompleteSd) --- the two things
wanted from an estimator that will be re-run every quarter. It is worth roughly
\pc{\nbThreeGainMean} of mean absolute error, not \pc{\nbThreeGainOne}. This is the one place in
the chapter where the Bayesian layer moves the estimate, and it moves it because the classical
toolbox had no third option to offer --- not because its arithmetic is better.

\subsection*{On the decision, the classical arm can play --- once the bar stops being rigged}

The only grade that finally matters is the action: \emph{price this region, or not?}
\Cref{sec:pe:euro}'s rule is $\Prob(\beta_r < -1) > 0.9$, which is a \textbf{one-sided 90\,\%}
requirement. Its like-for-like classical partner is therefore the \textbf{one-sided 90\,\% upper
confidence bound},
%
\begin{equation}
  \hat\beta_r \;+\; t_{0.90,\ \mathrm{df}_r}\, \widehat{se}_r \;<\; -1 ,
  \label{eq:pe:bound}
\end{equation}
%
the same 90, on the same side. The tempting alternative --- ``the two-sided 90\,\% confidence
interval lies entirely below $-1$'' --- looks natural and is \emph{wrong}: a two-sided 90\,\%
interval puts 5\,\% in each tail, so that test is secretly a one-sided \textbf{95\,\%} bar,
strictly stricter than what the posterior is being asked to clear. Scoring the classical arm there
would win the argument with a threshold instead of an estimator. \Cref{tab:pe:counts} scores it at
\cref{eq:pe:bound}, and reports both.

\input{build/tables/nb03_counts}

Matched honestly, no pooling acts on \nbThreeActNp{} regions (\nbThreeActNpCorrect{} correct,
\nbThreeActNpWrong{} wrong) against the posterior's \nbThreeActPp{} (\nbThreeActPpCorrect{}
correct, \nbThreeActPpWrong{} wrong). \textbf{The honest gap is \nbThreeGapHonest{} region.} At the
mismatched bar the same fits would act on only \nbThreeActNpStrict, and the gap would read
\nbThreeGapMismatched{} --- of which \nbThreeNFlip{} regions were never missed by the classical
\emph{estimator} at all. They were missed by the \emph{bar}: \nbThreeFlipRegions{} have one-sided
90\,\% bounds of \nbThreeFlipBounds, just past $-1$. The finding still points the same way, which is
why the chapter's conclusion survives. It simply stops being inflated.

Where the classical arm \emph{does} fail on this decision, fair scoring does not save it, and it
fails expensively. \textbf{Complete pooling's bound also clears $-1$, so it certifies all twelve
regions as elastic} --- including \nbThreeInelasticRegion, which the simulator planted as genuinely
inelastic at $\nbThreeInelasticBeta$ and which would consequently receive a price cut it cannot
support. That is a \emph{wrong} action, not a missed one, and it is the only wrong action anywhere
in \cref{tab:pe:counts}. No pooling, scored fairly, leaves \nbThreeActNpMissed{} of the elastic
regions on the table; the posterior leaves \nbThreeActPpMissed. On counts, that is a rounding
error --- and it is the right moment to stop pretending the win is about counts.

\subsection*{The dial exists. It just does not turn the knob the CMO needs}

It would be convenient to say that the classical arm has no threshold to tune. It has one:
$\alpha$. An analyst can build an 80\,\% bound, or test at any level she likes, and the roster of
priced regions re-sorts --- exactly as the posterior's roster re-sorts when the action bar moves
from 0.9 to 0.8. Anyone claiming the Bayesian arm is the only one with a dial has not looked.

The difference is \emph{what the two dials control}, and it survives the corrected scoreline.
Turning $\alpha$ tunes the \textbf{long-run error rate of the procedure}: at $\alpha = 0.1$, one in
ten bounds built this way lands on the wrong side of the truth, \emph{across hypothetical
repetitions of the whole study}. It cannot be set to a statement of risk appetite about
\textbf{\nbThreeShowcaseRegion{} in particular}, because ``the probability that
\nbThreeShowcaseRegion{} is elastic'' is not a quantity the procedure emits.
That region's bound either clears $-1$ or it does not (here it is \nbThreeShowcaseBound), and which
of those two things happened is not something the interval can tell anyone. The posterior's dial is
a different object entirely: for the same region, $\Prob(\beta_r < -1) = \nbThreeShowcaseP$ is a
number \emph{about that region}, in the only sample there is, and moving the bar to 0.8 changes
which regions the firm is willing to be wrong about, and how often. \textbf{Matching the \emph{level} of the two bars ---
which \cref{tab:pe:counts} does, honestly --- does not make them the same \emph{quantity}.}

And that is what pays for \cref{sec:pe:euro}'s third tier. Because each region arrives carrying a
number in $[0,1]$, the pricing rule can \textbf{route} it: act above 0.9, run a controlled price
test between 0.5 and 0.9, hold below. A confidence bound has no ``how sure are we'' to report --- a
region is in or it is out --- so the classical arm can build a two-tier policy and no more. The
posterior's \nbThreeActPpMissed{} misses are not a failure; they are a \emph{choice}, made against
a stated probability, and the regions it holds back are the ones \cref{sec:pe:euro} sends to a
price test rather than abandons.

\subsection*{The ledger, stated plainly}

What Bayes did \textbf{not} buy here. Not the fleet elasticity: \cref{sec:pe:classical} had it,
correctly, with an honest clustered interval, and the two agree to \nbThreeMuVsOls. Not a better
standard error: clustering already reported that the pooled model was misspecified, which is the
\emph{same message} the hierarchy acts on, merely delivered without a way to respond to it. Not
protection from endogeneity: \cref{sec:pe:wrong} shows the posterior exactly as biased as the OLS
slope when \cref{as:pe:exog} fails, and the fix is an instrument that \cref{sec:pe:classical}'s
toolkit fits with no prior in sight. And not, on a matched bar, a landslide on the decision count:
that margin is \nbThreeGapHonest{} region.

What it \emph{did} buy is two things, and both are real. The \textbf{third pooling regime} --- worth
\pc{\nbThreeGainMean} of mean absolute error on average and a materially tighter spread across
panels, with the weight between the extremes chosen by the data through $\tau_\beta$ rather than
asserted. And the \textbf{probability statement}: twelve honest per-region elasticities instead of
one, each carrying a $\Prob(\beta_r < -1)$ that a pricing rule can actually be written in. On the
euros, the edge over a fairly-scored classical arm is narrow. On the \emph{vocabulary}, the gap is
total --- and it is the vocabulary, not the euros, that lets \cref{sec:pe:euro} run a three-tier
\textsc{set} / \textsc{test} / \textsc{hold} policy instead of a coin flip.

% =====================================================================================
\section{The decision, in euros}
\label{sec:pe:euro}

The decision-maker's question was never ``what is $\beta$?''. It is ``\textbf{what price do I
set?}'' \Cref{eq:pe:lerner} answers it --- $P^\star = c\,\beta/(\beta+1)$ at a marginal cost of
$c = \eur{\nbThreeMarginalCost}$ --- but only for $\beta < -1$, and only if one is willing to plug
in a number one is not sure of.

\subsection*{Why the posterior is propagated draw by draw, and not plugged in at the mean}

$P^\star$ is a violently non-linear function of $\beta$ near the decision line: as
$\beta \to -1^{-}$ the denominator $\beta + 1 \to 0^{-}$ and $P^\star \to +\infty$. Posterior draws
sitting near $-1$ therefore produce wild optimal prices, and the mean of the resulting prices is
\emph{not} the price computed at the mean elasticity. That is Jensen's inequality, and here it is
not a technicality but the entire safety mechanism.

The chapter's own worked illustration: for the least confident region, \nbThreeWeakRegion, the
posterior mean elasticity is $\nbThreeWeakBeta$ and $\Prob(\beta < -1) = \nbThreeWeakP$. Plug that
mean into \cref{eq:pe:lerner} and it returns \eur{\nbThreeWeakPlug} --- against
\eur{\nbThreeStrongPstar} for the most confident region, \nbThreeStrongRegion. The estimator has
not gone mad; it has been asked a question it is not entitled to answer. So every quantity in this
section is computed \emph{per posterior draw} --- elasticity, optimal price, profit uplift, all
consistent within a draw --- and only then averaged.

\subsection*{The rule: three tiers, not two}

\begin{equation}
  \Prob(\beta_r < -1)\ \begin{cases}
    > 0.9 & \Rightarrow\ \textsc{set price} \quad\text{(move to $P^\star$, staged)},\\[2pt]
    \in (0.5,\ 0.9] & \Rightarrow\ \textsc{controlled test} \quad\text{(manufacture the variation)},\\[2pt]
    \le 0.5 & \Rightarrow\ \textsc{hold}.
  \end{cases}
  \label{eq:pe:rule}
\end{equation}
%
Applied to the twelve regions (\cref{fig:pe:decision}, \cref{tab:pe:pricing}), \cref{eq:pe:rule}
prices \nbThreeNSet{} regions, sends \nbThreeNTest{} to a controlled test, and leaves
\nbThreeNHold{} in the bottom \textsc{hold} tier. Across the priced regions, moving to $P^\star$ lifts profit by
\pc{\nbThreeFleetUplift} on average, with a 90\,\% credible interval of
$[\pc{\nbThreeFleetUpliftLo},\ \pc{\nbThreeFleetUpliftHi}]$ --- the uncertainty in the elasticity
carried all the way through to the euro, which is the entire point of having a posterior.

\input{build/floats/nb03_decision}
\input{build/tables/nb03_pricing}

Two subtleties in \cref{tab:pe:pricing} that a sharp audience will probe.

\emph{The uplift column is already discounted by our uncertainty.} It is a posterior mean over
\emph{every} draw, with the draws in which the region turns out inelastic contributing $0\,\%$
(because raising price does not pay there). For the confident \textsc{set price} rows, where
$\Prob(\beta_r < -1) \approx 1$, this is the same as ``the uplift once we act''. For
\nbThreeWeakRegion{} it is not: the expected uplift of \pc{\nbThreeWeakUplift} is diluted by the
draws in which $\beta \ge -1$ and the uplift is worth nothing --- a fraction
$1 - \nbThreeWeakP$ of the posterior --- and \emph{conditional} on that region truly being elastic
it is \pc{\nbThreeWeakUpliftCond}, a big and thoroughly unreliable number. Which is exactly why the
region goes to a test rather than onto the price list.

\emph{No price is published for a region that did not clear the bar.} The $P^\star$ column reads
``---'' for the routed regions. Publishing a price for a region one has just declared too uncertain
to price would invite precisely the mistake the routing exists to prevent.

\subsection*{Why the low-confidence regions get a test, seen rather than argued}

\input{build/floats/nb03_profit}

\Cref{fig:pe:profit} makes the case visually, and it is the figure to show a CMO. For the most
confident region the posterior draws agree on a gentle hill with a \emph{flat top}: even a sizeable
error in $\beta$ moves profit very little near $P^\star$, so acting is cheap. For the least
confident region the curves fan out and many keep climbing off the right edge --- the model is
begging for a price far outside the green band of prices ever charged, on a functional form
(\cref{as:pe:form}) that could only ever be tested \emph{inside} that band. Statistical risk and
extrapolation risk stack, and the region is routed to a test.

Note the sobering detail visible even in the confident panel: $P^\star$ typically sits
\emph{above} the observed price range, which runs \eur{\nbThreePriceObsLo}--\eur{\nbThreePriceObsHi}
against recommended prices of \eur{\nbThreePstarLo}--\eur{\nbThreePstarHi}. \textbf{The
recommendation is itself an extrapolation.} That is why the operational advice is a staged
rollout --- move a minority of stores first, watch a few weeks, re-estimate as the new prices
generate data --- rather than a fleet-wide overnight change, even where the posterior is sure. And
it is why the \textsc{controlled test} tier is a deliverable rather than a shrug: randomized
$\pm 10\,\%$ price cells for eight weeks deliberately manufacture the exogenous price variation
that \cref{as:pe:variation} feeds on.

\subsection*{The cost assumption --- the other number in the formula}

\input{build/tables/nb03_cost}

$P^\star = c\,\beta/(\beta+1)$ is \emph{linear} in marginal cost. A $25\,\%$ error in $c$ moves
every recommended price by $25\,\%$, with no posterior required and no sampler consulted. Cost is
usually handed to the analyst as a single accounting number, and it deserves at least the
scepticism the elasticity gets --- arguably more, because nothing in the data can check it.
\Cref{tab:pe:cost} sweeps $c$ from \eur{\nbThreeCostLo} to \eur{\nbThreeCostHi} around the
\eur{\nbThreeMarginalCost} baseline, and the result is a clean separation worth remembering:

\begin{itemize}[leftmargin=*]
  \item \textbf{The routing never changes.} The set-price roster stays at \nbThreeNSet{} regions
        across the whole sweep, because \cref{eq:pe:rule} keys off $\Prob(\beta_r < -1)$, in which
        $c$ does not appear.
  \item \textbf{Every price changes.} $P^\star$ scales exactly linearly, and the uplift moves with
        it.
\end{itemize}

So if finance restates the cost, the set-versus-test split survives untouched and the price list is
rewritten wholesale. Cost measurement is as decision-critical as the elasticity itself, and it is
usually audited less.

% =====================================================================================
\section{What can go wrong}
\label{sec:pe:wrong}

Partial pooling handles \emph{noise}. It does not handle \emph{bias}, and the bias that matters in
pricing has a name.

\subsection*{Endogeneity, and why the hierarchy is no protection}

Set $\kappa = \nbThreeKappaDemo$ in \cref{eq:pe:price}: the price now responds to the same
unobserved demand shock $u_{rt}$ that moves demand in \cref{eq:pe:demand}. \Cref{as:pe:exog} is
broken. The truth is unchanged --- the planted fleet elasticity on this panel is
$\nbThreeTrueFleetEndog$ --- and the full hierarchical model of \cref{sec:pe:bayes} is refit on it,
priors and all.

\input{build/tables/nb03_endog}

The result is the most important negative finding in this chapter. The hierarchical posterior
returns a fleet elasticity of $\nbThreeHierEndog$, biased toward zero by $\nbThreeObsBias$ --- and
\cref{eq:pe:ovb} predicted $\nbThreePredBias$ from the planted constants alone, so the bias is
\emph{computable}, not mysterious. More pointedly still: that posterior sits within
\textbf{\nbThreeHierVsOls} of naive OLS on the same raw panel ($\nbThreeOlsEndogRaw$). The
hierarchy, the priors, the sampler, the borrowed strength --- none of it moved the number.

\begin{remark}[The sentence to take away]
\textbf{Partial pooling is not a substitute for identification.} It shrinks noise; endogeneity is
bias. Every pooling regime in \cref{sec:pe:compare} --- complete, none, partial --- is wrong by the
same amount when \cref{as:pe:exog} fails, because they all read the same confounded correlation and
merely disagree about how to average it.
\end{remark}

\subsection*{It does not shade the number. It reverses the call}

\input{build/floats/nb03_kappa}

One value of $\kappa$ is an anecdote, so \cref{fig:pe:kappa} sweeps it from $0$ to
\nbThreeKappaMax{} and overlays \cref{eq:pe:ovb}'s prediction. Two things to watch, and both are
worse than a gentle drift.

The estimate does not degrade gracefully. Because the honest price variation is tiny
($\sigma_\nu = \nbThreeSigmaNu$), the bias term
$\kappa\sigma_u^2/(\kappa^2\sigma_u^2 + \sigma_\nu^2)$ peaks near $\kappa \approx \sigma_\nu$, and
over a whole stretch of the sweep the fitted demand curve slopes \emph{upward} --- the model
reporting that customers buy more when the price rises. Theory and refit agree there too, which is
the point of drawing the grey line.

And the recommendation flips almost immediately. The estimate crosses the $\beta = -1$ decision
line by $\kappa \approx \nbThreeKappaCross$. A hidden shock explaining only about
\pc{\nbThreeKappaCrossShare} of the residual variance in price is enough to turn ``set optimal
prices in the elastic regions'' into ``demand looks inelastic, hold everywhere'' --- fleet-wide,
with confidence, and with every diagnostic in \cref{sec:pe:valid} still passing.

\subsection*{The cure, and its price}

\input{build/floats/nb03_iv}

The fix is not a better estimator of the same correlation. It is an \textbf{instrument}: a variable
that moves price and affects demand \emph{only} through price. In pricing, the canonical one is a
\textbf{cost shifter} --- an input-cost shock, an exchange-rate move, a freight surcharge --- which
by construction pushes the price without appearing in the demand equation. Planting one here and
running two-stage least squares \citep{imbens1994} gives the picture in \cref{fig:pe:iv} and the
last two rows of \cref{tab:pe:endog}.

Naive OLS is \textbf{precisely wrong}: $\nbThreeBetaOls \pm \nbThreeSeOls$, which is
\nbThreeOlsSesOff{} standard errors from the planted $\nbThreeTrueFleetEndog$, on the inelastic side
of the decision line, with a tight interval $[\nbThreeOlsEndogLo,\ \nbThreeOlsEndogHi]$ around it.
And its most \emph{efficient} standard error is its most confident lie: under HC1 the same estimate
carries a standard error of \nbThreeOlsSeHc{} instead of \nbThreeSeOls, which puts the truth
\nbThreeOlsSesOffHc{} standard errors away. A competent classical analyst, making no mistakes, would
confidently hold prices fleet-wide.

2SLS is \textbf{imprecisely right}: $\nbThreeBetaIv \pm \nbThreeSeIv$, landing
\nbThreeIvSesOff{} standard errors from the truth, with a first-stage $F$ of \nbThreeIvF{} ---
comfortably past the conventional weak-instrument bar of ten \citep{stock2005} --- and an interval
$[\nbThreeIvLo,\ \nbThreeIvHi]$ that is \nbThreeIvWidthRatio$\times$ the width of the OLS one. It
keeps only the cost-driven slice of price variation, so the hidden shock drops out of the
estimate --- and most of the data drops out with it.

That asymmetry \emph{is} the instrumental-variables bargain: the cure costs variance. Real IV
estimates are noisier and need bigger samples, and a \emph{weak} instrument (first-stage
$F \lesssim 10$) is worse than none, because it amplifies small exclusion violations into large
bias. Managing that trade-off --- with honest IV uncertainty, weak-instrument diagnostics, and a
Bayesian wrapper --- is the job of Chapter~13, which also repeats it on real advertising data. This
chapter's job was to show why the instrument is not optional.

\subsection*{The remaining limits}

\begin{itemize}[leftmargin=*]
  \item \textbf{\Cref{as:pe:form}, functional form.} Constant elasticity survives the in-sample
        probe of \cref{sec:pe:valid} --- with the honest note that $\pm 5\,\%$ of price wiggle gives
        that test very little power. Every $P^\star$ sits at or beyond the edge of the observed
        band, where no in-sample check reaches. The staged rollout of \cref{sec:pe:euro} is the
        operational hedge, and the elasticity should be re-estimated as the new prices generate
        data.
  \item \textbf{\Cref{as:pe:sutva}, dynamics.} A lagged-price term captures stockpiling and
        reference-price effects; it is omitted here for clarity. On real scanner data, always try
        it. Pantry-loading masquerades as high short-run elasticity and vanishes at longer horizons,
        and a firm that prices on the short-run number will discount its way into a hole.
  \item \textbf{Levels versus logs.} The model is for $\log Q$. Anyone forecasting \emph{units} from
        it needs the lognormal correction $\E[Q \mid P] = e^{\mu + \sigma^{2}/2}$, not $e^{\mu}$.
        The profit \emph{uplift} of \cref{sec:pe:euro} is immune, being a ratio of two expected
        profits with the same residual variance, so the $e^{\sigma^2/2}$ factor cancels --- but the
        moment absolute euros or units are quoted, it bites.
  \item \textbf{Pooling is a modelling choice, and it can be pushed too far.} Too few regions, or a
        prior on $\tau_\beta$ that is too tight, and the model over-pools --- every region gets the
        fleet answer, which is complete pooling with extra steps. Inspect $\tau_\beta$'s posterior
        (here \nbThreeTauBeta, comfortably away from zero) rather than trusting the machinery to
        notice.
  \item \textbf{And the cluster count does not improve with the data.} There are
        \nbThreeNClusters{} regions and there will still be \nbThreeNClusters{} regions next year.
        The classical interval of \cref{tab:pe:covariance} is mildly optimistic for that reason and
        will remain so; a wild cluster bootstrap \citep{cameron2008} is the repair when the decision
        is close.
\end{itemize}

% =====================================================================================
\section{Takeaways}
\label{sec:pe:take}

\begin{enumerate}[leftmargin=*, itemsep=.45em]
  \item \textbf{The model was chosen for interpretability, not convenience.} In a log-log demand
        curve the coefficient on log-price \emph{is} the elasticity (\cref{eq:pe:elast}), and by
        Lerner's rule (\cref{eq:pe:lerner}) the elasticity \emph{is} the optimal markup. There is no
        back-transformation and no marginal effect to compute --- which is also why a mis-estimated
        $\beta$ goes straight into a price.

  \item \textbf{The decision line is $-1$, not zero.} Revenue rises with price above it and falls
        below it (\cref{eq:pe:revenue}), so the sign of the instruction sent to a region manager
        flips there. That is why the rule is $\Prob(\beta_r < -1) > 0.9$ and not ``is it
        significant?'': significance is a statement about zero, and the business does not care about
        zero.

  \item \textbf{Partial pooling wins --- by \pc{\nbThreeGainMean} on average, not the
        \pc{\nbThreeGainOne} one panel advertises.} Over \nbThreeNseed{} fresh panels the margin has
        a standard deviation of \nbThreeGainSd\,pp and partial pooling \emph{loses} on
        \nbThreeNLosses{} of them. Worse, \emph{which} classical arm is the one to beat is itself a
        one-seed artefact: the working panel says complete pooling, the mean of \nbThreeNseed{}
        panels says no pooling (\nbThreeMaeNopool{} against \nbThreeMaeComplete). A chapter that
        lectures the reader that one sample cannot separate bias from luck must obey its own rule.

  \item \textbf{The classical arm has a dial, and matching the level does not match the quantity.}
        Scored at a matched one-sided 90\,\% bound (\cref{eq:pe:bound}), no pooling acts on
        \nbThreeActNp{} regions against the posterior's \nbThreeActPp{} --- an honest gap of
        \nbThreeGapHonest{} region, not the \nbThreeGapMismatched{} a mismatched threshold would
        have sold. What survives the correction is not the count but the vocabulary: for
        \nbThreeShowcaseRegion, $\Prob(\beta_r < -1) = \nbThreeShowcaseP$ is a statement about
        \emph{that region}, while its one-sided bound of \nbThreeShowcaseBound{} is a statement about
        the \emph{procedure}. Only the first can be routed into a three-tier \textsc{set} /
        \textsc{test} / \textsc{hold} policy.

  \item \textbf{Complete pooling's failure is a \emph{wrong} action, not a missed one.} One verdict
        imposed on twelve regions certifies \nbThreeInelasticRegion{} --- planted genuinely
        inelastic at $\nbThreeInelasticBeta$ --- as elastic, and would cut its price. That is the
        only wrong action in \cref{tab:pe:counts}, and it is the cost of pretending regions are
        identical.

  \item \textbf{Hierarchy is not identification.} On an endogenous panel the Bayesian posterior's
        fleet elasticity sits within \nbThreeHierVsOls{} of naive OLS on the same data. The bias is
        \emph{predictable} --- \cref{eq:pe:ovb} calls $\nbThreePredBias$ and the refit delivers
        $\nbThreeObsBias$ --- and it is \emph{decisive}: a hidden shock explaining only
        \pc{\nbThreeKappaCrossShare} of the residual variance in price flips the pricing call
        fleet-wide. \textbf{Partial pooling shrinks noise. Endogeneity is bias. The fix is an
        instrument.}

  \item \textbf{The covariance estimator should be chosen, and then distrusted.} Clustering on
        region widened the standard error \nbThreeSeRatioIid$\times$ over iid, and the widening was
        itself the diagnosis --- the pooled model's own slope misspecification, not serial
        correlation. And with $G = \nbThreeNClusters$ clusters, even the clustered interval is
        \emph{downward}-biased \citep{cameron2015}. The best interval the classical toolbox has here
        is still the optimistic one, and saying so is part of the estimate.

  \item \textbf{Every recommended price is an extrapolation.} The observed band is
        \eur{\nbThreePriceObsLo}--\eur{\nbThreePriceObsHi}; the recommendations are
        \eur{\nbThreePstarLo}--\eur{\nbThreePstarHi}. Constant elasticity was only ever tested where
        the data lives. Stage the rollout, and re-estimate on what it generates.
\end{enumerate}

\notebooknote{%
  This chapter is \code{notebooks/03\_price\_elasticity.ipynb}. It runs in a fast teaching mode
  (\code{CMP\_FAST=1}, short chains, \nbThreeNseed{} panels reduced) and a full mode
  (\code{CMP\_FAST=0}); every number here comes from the full run, emitted by \code{cmp.report} and
  injected as a macro. The simulator is \code{cmp.dgp.price\_panel}, the classical estimators
  (\code{ols}, \code{no\_pooling}, \code{iv\_2sls}) are \code{cmp.classical}, the hierarchical model
  is PyMC, and the seeds are fixed in the notebook. The instrumental-variables treatment previewed
  in \cref{sec:pe:wrong} is Chapter~13.}
